\raggedright{\hspace{3mm}

\centerline{\bf \Huge 5 класс}

\problem{\textbf{1}} 
Расул Владимирович придумал систему перевода цифр в буквы и обратно, в которой каждой из десяти цифр соответствует одна уникальная буква. Расул Владимирович решил поделиться своей системой с Анджуром Аркадьевичем. Анджур Аркадьевич решил воспользоваться интересной системой и умножил в столбик два слова: ЭЛМШ с ЭТЛ, получив в произведении ЭТЛЭТЛ. Определите все возможные варианты чисел, скрывающихся за ЭЛМШ с ЭТЛ, и докажите, что других вариантов быть не может.

\problem{\textbf{2}} 
В ЭлМШ учатся Шиншиллы, которые всегда лгут, Лососи, которые всегда говорят правду, и Киви, которые могут сказать что угодно по своему усмотрению. Однажды 2025 учеников ЭлМШ построились \textit{черепахой} (квадрат из 45 строк и 45 столбцов). Затем каждый сказал: «В строке, в которой я стою, Шиншилл больше, чем остальных вместе взятых в этой строке, а в столбце меньше, чем остальных вместе взятых в этом столбце». Какое наименьшее и наибольшее количество Шиншилл может быть в данном построении \textit{черепахой}?

\problem{\textbf{3}} 
Можно ли разрезать квадрат $12\times12$ по линиям сетки на какое-то количество фигур первого вида и какое-то количество фигур второго вида так, чтобы фигур первого вида было в 4 раза больше, чем фигур второго, и периметры фигур обоих видов были равны?

\problem{\textbf{4}} 
Айна загадала два различных натуральных числа, произведение которых больше 14. Тагир загадал два различных натуральных числа, больших чем числа Айны. Тимофей загадал ещё два различных натуральных числа, больших чем числа Тагира, и произведение которых меньше 89. Найдите все возможные варианты этих шести чисел и докажите, что других вариантов быть не может (натуральные числа ‒ это 1, 2, 3, 4, 5, 6, 7, 8, 9, 10, 11, 12, 13 и так далее).

\problem{\textbf{5}} 
В олимпиаде участвуют 100 учеников (вам незнакомых), из которых 85 – \textit{УченикиЭТЛ} и 15 –\textit{Неученики}. Расул Владимирович знает всех в лицо и поэтому может за одну чокопайку определить любую пару учеников, то есть сказать, кто они: «оба УченикиЭТЛ», «оба Неученики», «оба разные». Можно ли, имея 64 чокопайки, найти всех \textit{УчениковЭТЛ}?}


